\label{sec:revisao-bibliografica}

Esta seção está dividida em duas partes. A fundamentação teórica (Seção~\ref{sec:fundamentacao}) apresenta os conceitos que sustentam a pesquisa: os SDKs de terceiros e o rastreamento em aplicativos móveis, as formas de detectar esses componentes e o marco brasileiro de proteção de dados. Os trabalhos correlatos (Seção~\ref{sec:correlatos}) reúnem os estudos sobre rastreamento e regulação, sobre o confronto entre o que os aplicativos declaram e o que fazem e sobre aplicativos de governo, além dos estudos com proposta de inclusão na revisão sistemática, cujo protocolo é descrito na Seção~\ref{sec:revisao}.

\subsecao{Fundamentação teórica}\label{sec:fundamentacao}

\subsubsecao{Bibliotecas de terceiros e rastreamento em aplicativos móveis}\label{sec:ref-rastreamento}

Aplicativos Android raramente são escritos do zero. Eles incorporam bibliotecas e SDKs de terceiros, que facilitam o desenvolvimento e oferecem funções como análise de uso, integração com redes sociais e monetização por anúncios \cite{razaghpanah2018}. Quando esses componentes coletam dados de comportamento e permitem identificar o usuário e acompanhá-lo em diferentes serviços digitais, fala-se em rastreamento por terceiros \cite{binns2018}. Como funcionam dentro do aplicativo, esses serviços são, em grande parte, invisíveis para o usuário e podem rastreá-lo sem consentimento, inclusive entre aplicativos diferentes \cite{razaghpanah2018}.

\subsubsecao{Detecção de bibliotecas e de rastreadores}\label{sec:ref-deteccao}

A presença de bibliotecas e os fluxos de dados que elas podem produzir são examinados por abordagens complementares. A análise estática examina o pacote de instalação (APK) sem executá-lo: lista as classes do código compilado e os textos que ele contém e os compara com assinaturas conhecidas. A análise dinâmica executa o aplicativo e observa o que ele de fato transmite, como fizeram Razaghpanah et~al.~\cite{razaghpanah2018}. A primeira mostra que o código está no pacote, mas não que é executado; a segunda mostra o envio, mas exige executar cada aplicativo.

Para a análise estática de bibliotecas, há várias ferramentas. Zhan et~al.~\cite{zhan2022} compararam as ferramentas públicas de detecção de bibliotecas de terceiros em Android quanto à eficácia, medida pela precisão e pela revocação, à identificação de versões, à eficiência, à resistência à ofuscação de código e à facilidade de uso. Nenhuma ferramenta foi a melhor em todos os critérios: o LibScout foi o mais eficaz; o LibRadar, o mais rápido e o mais fácil de usar; o ORLIS, o que melhor resistiu a diferentes ofuscadores; e o LibPecker, o de melhor desempenho geral nos testes com ofuscação. Como a ofuscação afeta cada ferramenta de modo diferente, ferramentas diferentes podem dar respostas diferentes para o mesmo aplicativo, e é esse efeito, aplicado a detectores de rastreadores, que a comparação entre detectores deste projeto mede.

Para rastreadores, uma base usada por ferramentas abertas é a do Exodus Privacy, que associa a cada rastreador conhecido uma assinatura de código, isto é, um padrão nos nomes das classes do aplicativo, e uma assinatura de rede, um padrão nos domínios que ele contata \cite{exodus2026}. Outras listas, como a da Disconnect, classificam domínios por finalidade, como publicidade, análise de uso e coleta de impressão digital do aparelho \cite{disconnect2026}. Ferramentas de análise de aplicativos, como o Mobile Security Framework (MobSF) \cite{mobsf2026}, e bibliotecas de extração, como o Androguard \cite{androguard2026}, permitem aplicar essas bases a um pacote, mas não são todas independentes entre si (Seção~\ref{sec:deteccao}). Combinar ferramentas exige, por isso, comparar unidades equivalentes e documentar as dependências entre bases, métodos de extração e critérios de decisão.

A origem dos pacotes também importa. A Google Play entrega a versão vigente, que é a que o cidadão instala, mas essa versão muda a cada atualização e pode ter a distribuição restrita por região. O AndroZoo é um acervo acadêmico com milhões de pacotes Android, no qual cada versão de um aplicativo é identificada pelo hash SHA-256 \cite{allix2016,alecci2024}, de modo que outro pesquisador pode obter exatamente o mesmo arquivo. A escolha entre as duas fontes é discutida na Seção~\ref{sec:pacotes}.

\subsubsecao{Proteção de dados pessoais e soberania de dados}\label{sec:ref-lgpd}

No Brasil, o tratamento de dados pessoais é regido pela LGPD \cite{brasil2018}. A lei define dado pessoal como a informação relacionada a pessoa natural identificada ou identificável e dado pessoal sensível como, entre outros, o dado referente à saúde e o dado genético ou biométrico, quando vinculado a uma pessoa natural (Art. 5º, I e II). O tratamento só pode ocorrer nas hipóteses previstas em lei (Art. 7º), e as hipóteses para dados sensíveis são mais restritas (Art. 11). O tratamento feito pelo poder público deve atender à sua finalidade pública, na persecução do interesse público (Art. 23). A transferência internacional, definida como a transferência de dados pessoais para país estrangeiro ou organismo internacional (Art. 5º, XV), só é permitida nos casos que a lei enumera (Art. 33).

A transparência é um dos princípios da LGPD: o titular deve receber informações claras, precisas e facilmente acessíveis sobre o tratamento e os respectivos agentes (Art. 6º, VI), inclusive sobre o uso compartilhado de seus dados e a finalidade desse compartilhamento (Art. 9º, V). No poder público, a lei exige ainda que os órgãos informem, de forma clara e atualizada e em veículos de fácil acesso, a previsão legal, a finalidade, os procedimentos e as práticas do tratamento (Art. 23, I) \cite{brasil2018}. Com base nesses dispositivos, este projeto considera que a coleta feita por um SDK de terceiro deve estar coberta pelas informações prestadas ao cidadão, ainda que sem o nome do componente.

A Agência Nacional de Proteção de Dados (ANPD), criada como Autoridade Nacional de Proteção de Dados, passou a ter natureza de autarquia especial em 2022 \cite{brasil2022}, e a Lei nº 15.352/2026 tornou-a agência reguladora, vinculada ao Ministério da Justiça e Segurança Pública \cite{brasil2026}. As orientações publicadas antes dessa mudança conservam a autoria institucional da época. O guia da ANPD para o poder público orienta a análise de bases legais, finalidades e compartilhamentos \cite{anpd2023}. O guia sobre cookies declara que suas orientações valem, de forma geral, para tecnologias semelhantes de rastreamento em dispositivos móveis \cite[p.~6]{anpd2022} e pede que a finalidade de cada categoria de cookie seja específica e informada ao titular \cite[p.~14]{anpd2022}, o que apoia o confronto entre os SDKs embarcados e o que cada aplicativo declara. Para a dimensão internacional, a Resolução CD/ANPD nº 19/2024 distingue transferência e coleta internacional e disciplina os mecanismos de transferência \cite{anpd2024}.

Como o rastreamento se concentra em poucas empresas e opera em várias jurisdições \cite{binns2018,kollnig2021before}, a presença de fornecedores estrangeiros torna relevante investigar a governança dessa dependência, sem presumir que o órgão desconheça a decisão técnica ou que todo SDK envie dados ao exterior. Pohle e Thiel~\cite{pohle2020} discutem a soberania digital em termos das capacidades de decisão de Estados, organizações e indivíduos e destacam que o conceito tem sentidos distintos. Neste projeto, soberania de dados é um recorte operacional dessa discussão: a capacidade de o Estado saber e decidir para onde vão os dados gerados no uso de seus serviços. A presença de SDKs de empresas estrangeiras já indica dependência de fornecedores estrangeiros; afirmar que há dados no exterior exige observar o tráfego, e caracterizar uma transferência internacional exige ainda identificar os agentes e a operação (Seção~\ref{sec:trafego}).

\subsecao{Trabalhos correlatos}\label{sec:correlatos}

Os trabalhos correlatos estão organizados em quatro grupos: os estudos sobre rastreamento em aplicativos e o efeito da regulação, os que confrontam o que o aplicativo declara com o que ele faz, os que examinam aplicativos de governo e, por fim, a síntese dos estudos com proposta de inclusão na revisão sistemática, que compara esses grupos e aponta a lacuna que este projeto procura preencher.

\subsubsecao{Rastreamento em aplicativos e regulação da proteção de dados}\label{sec:ref-regulacao}

Binns et~al.~\cite{binns2018} mediram, por análise estática do código dos pacotes, a presença de rastreadores em 959 mil aplicativos das lojas da Google Play dos Estados Unidos e do Reino Unido. A maioria dos aplicativos continha rastreamento de terceiros, e a distribuição era concentrada: poucos rastreadores, como os do Google, apareciam em grande parte dos aplicativos, enquanto muitos outros apareciam em poucos, formando uma cauda longa. A intensidade variava por categoria, com aplicativos de notícias e infantis entre os que mais embarcavam rastreadores. O estudo mostrou ainda que o rastreamento é transnacional, com muitos rastreadores operando em jurisdições fora da União Europeia.

Razaghpanah et~al.~\cite{razaghpanah2018} partiram do tráfego real de usuários, coletado por um aplicativo de monitoramento instalado nos próprios aparelhos, e desenvolveram métodos para detectar serviços de publicidade e rastreamento no nível da rede. Identificaram 2.121 desses serviços, dos quais 233 não constavam das listas de bloqueio mais usadas. Ao analisar as políticas de privacidade dos maiores fornecedores, concluíram que o compartilhamento dos dados coletados com subsidiárias e afiliadas é a norma. O resultado expõe um limite das medições baseadas em listas: elas só encontram o que já conhecem.

Comparações entre períodos anteriores e posteriores a uma mudança regulatória descrevem tendências, mas não isolam, por si sós, o efeito de uma lei. Kollnig et~al.~\cite{kollnig2021before} compararam quase dois milhões de aplicativos Android, cerca de um milhão de 2017 e um milhão de 2020, ou seja, de antes e de depois da entrada em vigor do Regulamento Geral de Proteção de Dados europeu (GDPR), em 2018. Procurando rastreadores no código dos pacotes, observaram pouca mudança na presença de rastreadores e a persistência da concentração da capacidade de rastrear em poucas grandes empresas.

Kollnig et~al.~\cite{kollnig2021fait} mostraram que a legislação da União Europeia e do Reino Unido exige consentimento para o rastreamento por terceiros, mas que a maioria dos aplicativos analisados rastreia e poucos obtêm o consentimento antes de fazê-lo. As próprias bibliotecas de rastreamento atribuem ao desenvolvedor do aplicativo a tarefa de obter esse consentimento, sem oferecer meios robustos para isso. Nguyen et~al.~\cite{nguyen2022}, por sua vez, estudaram os avisos de consentimento em 239.381 aplicativos: 30.160 compartilhavam dados pessoais com terceiros sem sequer tentar obter consentimento e, entre os 13.082 que exibiam aviso, 2.688 (20,54\%) violavam ao menos um requisito do GDPR, por exemplo ao continuar transmitindo dados depois de o usuário recusar.

\subsubsecao{O que o aplicativo declara e o que ele faz}\label{sec:ref-politicas}

A política de privacidade descreve práticas que nem sempre são fáceis de identificar no código. Zimmeck et~al.~\cite{zimmeck2019} automatizaram o confronto entre o que as políticas declaram e o que o código dos aplicativos indica com o sistema MAPS, aplicado a aproximadamente um milhão de aplicativos.

O PoliCheck, de Andow et~al.~\cite{andow2020}, mostra que não basta saber a categoria do dado: também é preciso identificar quem o recebe. No conjunto de 13.796 aplicativos estudados, os autores verificaram que até 42,4\% deles declaravam de forma incorreta ou omitiam seus fluxos de dados sensíveis. A hipótese deste projeto segue essa linha de pesquisa, mas volta a atenção para os SDKs e para as declarações disponíveis na Google Play. Por isso, cada registro distinguirá o componente, o destinatário, a categoria de dado e a finalidade declarada.

\subsubsecao{Aplicativos de governo}\label{sec:ref-governo}

Há pesquisas específicas sobre aplicativos de governo, embora nenhuma resolva a questão proposta para o Brasil. Samarasinghe et~al.~\cite{samarasinghe2022} lembram que, diferentemente dos sites comerciais, os serviços de governo são financiados com dinheiro público, tratam dados sensíveis e raramente aparecem nas listas dos sites mais acessados, de onde costumam sair as amostras dos estudos de rastreamento. Os autores localizaram aplicativos pelos links para a Google Play presentes em sites de governo e examinaram 1.166 aplicativos Android de 71 países: 37\% continham rastreadores do Google, e 27\% enviavam informações sensíveis, como identificadores do usuário ou do dispositivo, senhas e chaves de acesso a interfaces de programação (API), a terceiros ou por conexões sem criptografia. Parte dos aplicativos não pôde ser baixada por restrição geográfica, e os resultados não são apresentados separadamente para o Brasil nem por esfera de governo.

Carlsson et~al.~\cite{carlsson2022} compararam o tráfego de rede de 32 aplicativos do setor público finlandês com as respectivas políticas de privacidade. Encontraram divergências entre o que as políticas diziam e os dados enviados a serviços de análise de terceiros, incluindo, em alguns casos, dados enviados para fora da Europa quando a política afirmava o contrário. É o estudo mais próximo do que este projeto propõe, embora trate de poucos aplicativos e de outro país.

Em estudo sobre aplicativos da COVID-19, Serrano et~al.~\cite{serrano2024} identificaram rastreadores em 402 de 595 aplicativos por meio das assinaturas do Exodus Privacy e examinaram o código de bibliotecas para caracterizar a coleta e o envio de dados, com o apoio da ferramenta SAPITO, apresentada no próprio artigo. Os pacotes vieram da Google Play, do AndroZoo, de um conjunto publicado no Zenodo e de repositórios alternativos, o que evidencia a necessidade de registrar a origem e a versão de cada pacote.

Esses resultados ajudam a formular a pesquisa, mas deixam três questões em aberto no recorte adotado aqui: a medição do conjunto dos aplicativos governamentais brasileiros, por esfera de governo; o confronto entre cada SDK detectado e o que a política de privacidade e a seção ``Segurança dos dados'' declaram; e a concordância entre detectores aplicados às mesmas versões dos pacotes, que Samarasinghe et~al.~\cite{samarasinghe2022} não relatam. Para o poder público, sujeito à LGPD e às próprias declarações, a divergência entre o declarado e o encontrado é um sinal de alerta por si só, sem depender de uma comparação com o mercado.

\subsubsecao{Estudos com proposta de inclusão na revisão sistemática}\label{sec:ref-sintese}

A revisão sistemática (Seção~\ref{sec:revisao}) encaminhou 148 registros à leitura completa. Até 26 de novembro de 2026, 31 deles receberam proposta de inclusão, e 8 dessas propostas já foram confirmadas com a leitura integral do texto; as demais serão relidas na íntegra junto com os registros pendentes. A lista é, portanto, preliminar: as propostas ainda não confirmadas e os registros em aberto podem alterá-la. A Tabela~\ref{tab:sintese} apresenta os 31 estudos pela questão da revisão a que respondem.

\begin{table}[htbp]
\centering
\small
\caption{Estudos com proposta de inclusão por questão da revisão}\label{tab:sintese}
\begin{tabularx}{\textwidth}{|>{\raggedright\arraybackslash}p{3.6cm}|>{\raggedright\arraybackslash}X|}
\hline
\textbf{Questão} & \textbf{Estudos} \\
\hline
QR1. Rastreadores e envio de dados a terceiros (11) & Razaghpanah et~al.~\cite{razaghpanah2018}; Binns et~al.~\cite{binns2018b}; Binns et~al.~\cite{binns2018}; Hu et~al.~\cite{hu2019}; Kollnig et~al.~\cite{kollnig2021fait}; Kollnig et~al.~\cite{kollnig2021before}; Nguyen et~al.~\cite{nguyen2022}; Ahasanuzzaman et~al.~\cite{ahasanuzzaman2022}; Dong et~al.~\cite{dong2024}; Heid et~al.~\cite{heid2026}; Pybus et~al.~\cite{pybus2026} \\ \hline
QR2. Confronto entre o declarado e o que o aplicativo faz (8) & Kununka et~al.~\cite{kununka2018}; Andow et~al.~\cite{andow2020}; Verderame et~al.~\cite{verderame2020}; Zimmeck et~al.~\cite{zimmeck2021}; Zhao et~al.~\cite{zhao2023}; Xiao et~al.~\cite{xiao2023}; Tu et~al.~\cite{tu2024}; Baalous et~al.~\cite{baalous2025} \\ \hline
QR3. Aplicativos de governo, inclusive os da COVID-19 (10) & Samarasinghe et~al.~\cite{samarasinghe2022}; Carlsson et~al.~\cite{carlsson2022}; Lemos et~al.~\cite{lemos2022}; Jiang e Zheng~\cite{jiang2023}; Kouliaridis et~al.~\cite{kouliaridis2021}; Samhi et~al.~\cite{samhi2021}; Hatamian et~al.~\cite{hatamian2021}; Flood et~al.~\cite{flood2022}; Jordanoski et~al.~\cite{jordanoski2023}; Serrano et~al.~\cite{serrano2024} \\ \hline
QR4. Brasil, LGPD e transferência internacional (2) & Reis et~al.~\cite{reis2023}; Sterling et~al.~\cite{sterling2024} \\ \hline
\end{tabularx}
\fonte{Elaborado pelo autor (2026).}
\nota{QR: questão de revisão (Seção~\ref{sec:revisao}). Entre parênteses, o número de estudos com proposta de inclusão em cada questão.}
\end{table}

Cada grupo da tabela sustenta uma escolha do projeto. Os estudos de presença de rastreadores (QR1) separam o que o código contém do que o aplicativo envia: a busca no código mostra presença, mas não prova que cada domínio encontrado recebeu dados \cite{binns2018,kollnig2021before}, o que exige observar o tráfego \cite{razaghpanah2018,kollnig2021fait,nguyen2022}. Por isso, a hipótese trata da declaração dos SDKs presentes, e o envio será observado em uma subamostra.

Os estudos que confrontam o declarado com o realizado (QR2) usam métodos que não são equivalentes: o PoliCheck \cite{andow2020} parte de fluxos observados durante a execução, e outros trabalhos, de permissões e chamadas de API presentes no código \cite{tu2024}. O grupo mostra também que as práticas das bibliotecas de terceiros precisam ser analisadas à parte das do aplicativo \cite{zhao2023,verderame2020}, o que justifica tomar o par aplicativo--SDK como unidade e registrar separadamente componente, destinatário e finalidade. Baalous et~al.~\cite{baalous2025} já confrontam indicadores de coleta encontrados no código com a seção ``Segurança dos dados''. Como essa seção declara categorias de dados e finalidades, e não a lista de SDKs, a classificação deste projeto distingue a menção nominal, a declaração por categoria e a ausência de declaração compatível (Seção~\ref{sec:declaracoes}).

Os estudos sobre governo e transferência internacional (QR3 e QR4) mostram que aplicativos públicos já foram examinados em outros países, sobretudo os da COVID-19 \cite{kouliaridis2021,samhi2021,hatamian2021,flood2022,jordanoski2023,serrano2024}, e que o destino dos dados já foi tratado nas políticas de aplicativos públicos de saúde \cite{jiang2023} e no armazenamento no exterior dos dados de um aplicativo oficial \cite{sterling2024}. No Brasil, Lemos et~al.~\cite{lemos2022} analisaram o aplicativo Monitora Covid-19, considerando rastreadores, documentação e proteção de dados, e Reis et~al.~\cite{reis2023} analisaram políticas de privacidade de aplicativos à luz da LGPD, sem examinar o comportamento técnico dos aplicativos. Esses trabalhos apoiam a soberania de dados como eixo da discussão (Seção~\ref{sec:ref-lgpd}), mas nenhum trata do conjunto governamental brasileiro.

Assim, nenhum dos 31 estudos mede a presença de SDKs no conjunto dos aplicativos de governo brasileiros confrontando-a com as declarações de cada aplicativo, nem a concordância entre detectores aplicados às mesmas versões dos pacotes. A contribuição esperada responde a essas lacunas em três frentes: a análise dos 376 aplicativos oficiais identificados, nas esferas federal, estadual e municipal; a classificação de cada par aplicativo--SDK quanto ao que está declarado, que permite testar a hipótese; e a medida de quanto o resultado depende do detector escolhido.
